\documentclass[11pt]{article}

\usepackage{geometry}
\usepackage{amsmath}
\usepackage{amssymb}
\usepackage{graphicx}
\usepackage{hyperref}
\usepackage{cite}
\usepackage{listings}
\usepackage{xcolor}
\usepackage{algorithm}
\usepackage{algpseudocode}
\geometry{margin=1in}



\title{Knapsack distribuito: implementazione, comparazione e analisi}
\author{Matteo Ielacqua}

\lstset{
    language=Python,
    basicstyle=\ttfamily\small,
    keywordstyle=\color{blue},
    commentstyle=\color{gray},
    stringstyle=\color{red},
    breaklines=true,
    showstringspaces=false,
    numbers=left,
    numberstyle=\tiny
}

\begin{document}

\maketitle

\section{Introduzione}  
L'algoritmo dello zaino è un noto problema di ottimizzazione. L'idea è che si ha a disposizione uno zaino 
e diversi elementi ad ognuno dei quali è assegnato un peso e un profitto. L'obiettivo è di massimizzare il profitto mantenendo il peso
totale degli elementi entro un limite prefissato. Nella versione detta 0/1, ogni elemento può essere preso una sola volta,
il problema si pone quindi come  $$ max \sum v_ix_i $$
dove $v_i$ è il profitto dell'oggetto i-esimo e $x_i$ rappresenta la scelta (1 se scelto 0 altrimenti), soggetto a 
$$ \sum w_ix_i \leq W \land x_i \in \{0,1\} $$ 
\\
Le versioni più complete di questo problema prevedono che si possa scegliere lo stesso elemento più volte, sia con che senza limite superiore a queste scelte
(rispettivamente Bounded o Unbounded). Esistono anche versioni più complesse, come lo zaino multidimensionale, dove non esiste un solo vincolo 
(come nel caso più semplice il peso) ma esistono diversi vincoli da tenere in considerazione durante il calcolo. Questa classe di problemi gioca 
un ruolo chiave in diversi campi molto delicati ad esempio nelle operazioni logistiche o nello scheduling di task satellitari, in cui infatti 
stanno prendendo piede diverse soluzioni basate su euristiche.

\subsection{Soluzioni per knapsack 0/1}
Il problema è NP-completo, data una soluzione non esiste quindi un algoritmo che possa verificare se è ottima in tempo polinomiale.
Esiste tuttavia una soluzione che lo risolve in tempo pseudo-polinomiale basata sulla programmazione dinamica. 
La prima soluzione più semplice e meno ottimale ha complessità temporale $O(2^n)$ ed è basata su un albero decisionale esplorato ricorsivamente
\begin{algorithm}[H]
\caption{Knapsack 0/1 Ricorsivo}\label{alg:knapsack-rec}
\begin{algorithmic}
\Require $profits$, $weights$, $max\_weight$, $index$
\Ensure Profitto massimo ottenibile
\If{$index = 0$ \textbf{or} $max\_weight = 0$}
    \State \Return $0$
\EndIf
\State $pick \gets 0$
\If{$weights[index-1] \leq max\_weight$}
    \State $pick \gets profits[index-1] + \Call{Knapsack}{max\_weight - weights[index-1],\, profits,\, weights,\, index-1}$
\EndIf
\State $not\_pick \gets \Call{Knapsack}{max\_weight,\, profits,\, weights,\, index-1}$
\State \Return $\max(pick,\, not\_pick)$
\end{algorithmic}
\end{algorithm}


Le altre due soluzioni sono basate sulla programmazione dinamica. 
la prima usa un approccio top-down (memoization) proponendo sostanzialmente una naturale evoluzione dell'approccio precedente, 
ha una complessità temporale e spaziale $O(n*W)$

\begin{algorithm}[H]
\caption{Knapsack 0/1 Top-Down (memoization)}\label{alg:knapsack-memo}
\begin{algorithmic}
\Require $profits$, $weights$, $max\_weight$, $index$, $table$
\Ensure Profitto massimo ottenibile
\If{$index = 0$ \textbf{or} $max\_weight = 0$}
    \State \Return $0$
\EndIf
\If{$table[index][max\_weight] \neq -1$}
    \State \Return $table[index][max\_weight]$
\EndIf
\State $pick \gets 0$
\If{$weights[index-1] \leq max\_weight$}
    \State $pick \gets profits[index-1] + \Call{Knapsack}{max\_weight - weights[index-1],\, profits,\, weights,\, index-1,\, table}$
\EndIf
\State $not\_pick \gets \Call{Knapsack}{max\_weight,\, profits,\, weights,\, index-1,\, table}$
\State $table[index][max\_weight] \gets \max(pick,\, not\_pick)$
\State \Return $table[index][max\_weight]$
\end{algorithmic}
\end{algorithm}

la seconda è una versione bottom-up (tabulation) con la medesima complessità ma che evita la ricorsione
\begin{algorithm}[H]
\caption{Knapsack 0/1 Bottom-Up (tabulation)}\label{alg:knapsack-tab}
\begin{algorithmic}
\Require $profits$, $weights$, $max\_weight$, $len$
\Ensure Profitto massimo ottenibile
\State $table \gets \text{Matrix}(len+1,\ max\_weight+1)$
\For{$i \gets 0$ \textbf{to} $len$}
    \For{$j \gets 0$ \textbf{to} $max\_weight$}
        \If{$i = 0$ \textbf{or} $j = 0$}
            \State $table[i][j] \gets 0$
        \Else
            \State $pick \gets 0$
            \If{$weights[i-1] \leq j$}
                \State $pick \gets profits[i-1] + table[i-1][j - weights[i-1]]$
            \EndIf
            \State $not\_pick \gets table[i-1][j]$
            \State $table[i][j] \gets \max(pick,\ not\_pick)$
        \EndIf
    \EndFor
\EndFor
\State \Return $table[len][max\_weight]$
\end{algorithmic}
\end{algorithm}

Una soluzione più raffinata di questa prevede di usare solo l'ultima riga calcolata, in modo da ridurre la complessità spaziale a $O(W)$

\begin{algorithm}[H]
\caption{Knapsack 0/1 Bottom-Up (spazio ottimizzato)}\label{alg:knapsack-opt}
\begin{algorithmic}
\Require $profits$, $weights$, $max\_weight$, $len$
\Ensure Profitto massimo ottenibile
\State $table \gets \text{array}[0 \dots max\_weight]$ inizializzato a 0
\For{$i \gets 1$ \textbf{to} $len$}
    \For{$j \gets max\_weight$ \textbf{downto} $weights[i-1]$}
        \State $table[j] \gets \max(table[j],\ table[j - weights[i-1]] + profits[i-1])$
    \EndFor
\EndFor
\State \Return $table[max\_weight]$
\end{algorithmic}
\end{algorithm}

Lo svantaggio principale delle soluzioni basate su programmazione dinamica è che non lavorano bene su casi in cui i pesi non sono interi
(o non possiedono una differenza prefissata tra di loro). Lo stesso vale per l'approccio naive descritto nella prima soluzione.

\subsection{Soluzione Greedy}
La soluzione Greedy ha un indubbio vantaggio rispetto a tutte le altre, può scegliere anche frazioni di un oggetto, 
l'idea è quella di prendere gli oggetti basandosi sul loro rapporto peso/valore, prendendoli in ordine decrescente di rapporto,
finchè lo zaino non è pieno. Se un oggetto non è possibile prenderlo per intero allora si può prendere una sua frazione a seconda 
della capacità rimanente. 

\begin{algorithm}[H]
\caption{Fractional Knapsack (Greedy)}\label{alg:fknapsack}
\begin{algorithmic}
\Require $profits$, $weights$, $max\_weight$, $len$
\Ensure Profitto massimo ottenibile
\State $items \gets \text{sort}(\text{zip}(profits, weights), \text{key} = \frac{profit}{weight} \text{ decrescente})$
\State $profit \gets 0.0$
\State $current\_capacity \gets max\_weight$
\For{$i \gets 0$ \textbf{to} $len - 1$}
    \If{$items[i].weight \leq current\_capacity$}
        \State $profit \gets profit + items[i].profit$
        \State $current\_capacity \gets current\_capacity - items[i].weight$
    \Else
        \State $profit \gets profit + \frac{items[i].profit}{items[i].weight} \times current\_capacity$
        \State \textbf{break}
    \EndIf
\EndFor
\State \Return $profit$
\end{algorithmic}
\end{algorithm}


% \subsection{Branch and bound}
% Il Knapsack Branch and Bound è l'algoritmo più adottato nei solver industriali. Presenta una complessità temporale $O(2^n)$ ma mantiene una 
% complessità spaziale di $O(N)$ e può lavorare con frazioni di peso. 

% \begin{verbatim}

% define Node = {level, profit, weight}
% define Item = {weight,profit}

% define bound(upper_node,len,max_weight,nodes)
% {
%     if(upper_node.weight >= W)
%     {
%         return 0
%     }
%     profit_bound = upper_node.profit 
%     j = upper_node.level + 1
%     total_weight = upper_node.weight
%     while(j < n and total_weight + nodes[j].weight <= max_weight)
%     {
%         total_weight += nodes[j].weight
%         profit_bound += nodes[j].profit
%         j += 1
%     }
%     if(j<n)
%     {
%         profit_bound += (max_weight - total_weight)*nodes[j].value/nodes[j].weight
%     }
%     return profit_bound
% }

% define Knapsack(max_weight, nodes, len)
% {
%     nodes = nodes.sort((a,b)=> (b.profit/b.weight) - (a.profit/a.weight))
%     priority_queue = PriorityQueue()
%     dummy = Node{-1,0,0}
%     priority_queue.put(dummy)
%     max_profit = 0
%     while(not priority_queue.empty())
%     {
%         node = priority_queue.get()
%         if(node.level == -1)
%         {
%             v = Node{0,0,0}
%         }
%         else if(node.level == len -1)
%         {
%             continue
%         }
%         else
%         {
%             //Nodo senza considerare il prossimo item 
%             v = Node{node.level + 1, node.profit, node.weight}
%         }
%         v.weight += nodes[v.level].weight
%         v.profit += nodes[v.level].value 
        
%         if(v.weight <= max_weight and v.profit > max_profit)
%         {
%             max_profit = v.profit
%         }
%         v_bound = bound(v,len,max_weight,nodes)


%         if(v_bound > max_profit)
%         {
%             priority_queue.put(v)
%         }

%         //Nodo considerando il prossimo item ma senza aggiungerlo allo zaino 

%         v = Node{node.level+1, node.profit, node.weight}
%         v_bound = bound(v,len,max_weight,nodes)
%         if(v_bound > max_profit)
%         {
%             priority_queue.put(v)
%         }
%     }
%     return max_profit
% }
% \end{verbatim}

\section{Soluzioni parallele al Knapsack 0/1}

Le soluzioni a questo problema sono le più disparate e negli anni ne sono state scritte molte. In ambito aereospaziale il knapsack è usato sopratutto 
per lo scheduling di task operativi satellitari \cite{surveymethods} in particolare per l'osservazione terrestre \cite{exactmethodphoto}.
Un interessante modello di esecuzione è denominato COPA \cite{knapsackCopaSolution} e sviluppa un algoritmo parallelo cost optimal per 
risolvere il problema sia su CPU con OpenMP che su GPU con CUDA, l'idea dietro quest'algoritmo è la seguente: 1) Si prende 
l'array di oggetti V e lo si divide in 2 insiemi disgiunti. 2) Per il primo insieme si generano tutte le $N=2^{n/2}$ possibili 
combinazioni $a_i$, si calcola di ogni sottoinsieme $a_i$ peso e profitto $w_i,p_i$ e si aggiunge questa tripletta in una lista 
in ordine crescente. 3) Si esegue lo stesso con la seconda lista per ogni sottoinsieme $b_i$ e le triplette si mettono in una lista
di ordine non crescente. In questo modo si è generata la lista A, per generare la lista B si riproducono gli step precedenti 
semplicemente inserendo le triplette in una lista di ordine non crescente. A questo punto si può passare alla fase di ricerca
che prenderà in input le due liste ordinate.

\begin{algorithm}[H]
\caption{Algoritmo di ricerca}\label{alg:cap}
\begin{algorithmic}
\Require Due liste ordinate $A$ e $B$
\Ensure Soluzione finale $Bestvalue$ e $X$
\State Inizializza $MaxB_N \gets b_N.p$, $L_N \gets N$, $Bestvalue \gets 0$, $X_1 \gets (0, 0)$
\For{$i \gets N-1$ \textbf{downto} $1$}
    \If{$b_i.p > MaxB_{i+1}$}
        \State $MaxB_i \gets b_i.p$ \textbf{and} $L_i \gets i$
    \Else
        \State $MaxB_i \gets MaxB_{i+1}$ \textbf{and} $L_i \gets L_{i+1}$
    \EndIf
\EndFor
\State $i \gets 1$ \textbf{and} $j \gets 1$
\While{$i \leq N$ \textbf{and} $j \leq N$}
    \If{$a_i.w + b_j.w > c$}
        \State $j \gets j + 1$ \textbf{and continue}
    \EndIf
    \If{$a_i.p + MaxB_j > Bestvalue$}
        \State $Bestvalue \gets a_i.p + MaxB_j$ \textbf{and} $X_1 \gets (a_i,\, b_{L_j})$
    \EndIf
    \State $i \gets i + 1$
\EndWhile
\State Converti le due componenti decimali di $X_1$ in due numeri binari: $X_1 \gets (binary_1,\, binary_2)$
\State $X \gets \text{strcat}(binary_1,\, binary_2)$ \Comment{Concatena i due binari}
\end{algorithmic}
\end{algorithm}

Questo modello di esecuzione si può vedere subito che si presta bene ad essere parallelizzato, sia nella parte di generazione 
delle liste sia in quella poi di ricerca dell'ottimo.  

\subsection{COPA: Cost Optimal Parallel Algorithm}
L'algoritmo parallelo è stato pensato per essere risolto su un modello EREW (Exclusive Read, Exclusive Write) PRAM con shared memory e assume k processori accessibili. 
L'algoritmo sfrutta prima un'algoritmo di generazione parallela delle liste A e B

\begin{algorithm}[H]
\caption{Algoritmo parallelo di generazione}\label{alg:parallel_gen}
\begin{algorithmic}
\Require Elementi $(v_i, v_i.w, v_i.p)$, $i = 1, 2, \ldots, n/2$
\Ensure Lista ordinata $A_{n/2}$ (i.e. $A$)
\State  $A_0 \gets [(0, 0, 0)]$ e $IC \gets 1$
\For{$i \gets 0$ \textbf{to} $n/2 - 1$}
    \State $v_{i+1} \gets IC$
    \For{\textbf{all} $P_j$, $1 \leq j \leq k$, \textbf{in parallel}}
        \State Genera nuova lista $A'_i$ aggiungendo la tripletta $(v_{i+1}, v_{i+1}.w, v_{i+1}.p)$ a tutti gli elementi della lista ordinata $A_i$
        \State Chiama l'algoritmo di merge parallelo per unire $A_i$ e $A'_i$ in ordine non decrescente di peso, producendo la lista ordinata $A_{i+1}$
    \EndFor
    \State $IC \gets IC + 1$
\EndFor
\end{algorithmic}
\end{algorithm}

l'algoritmo di merging parallelo si articola nel modo seguente: prende in input due vettori ordinati $U = (u_1, u_2, \ldots, u_m)$ e $V = (v_1, v_2, \ldots, v_m)$ e $k$ processori $P_1, P_2, \ldots, P_k$, con $1 \leq k \leq 2m$ potenza di 2, e produce un vettore ordinato di lunghezza $2m$ ottenuto dall'unione di $U$ e $V$.
\begin{enumerate}
    \item \textbf{Step 1:} I $k$ processori partizionano $U$ e $V$ in parallelo, ciascuno in $k$ sottovettori (che possono essere vuoti) $(U_1, U_2, \ldots, U_k)$ e $(V_1, V_2, \ldots, V_k)$ tali che: $|U_i| + |V_i| = 2m/k$ per $1 \leq i \leq k$, e i pesi di tutti gli elementi in $U_i$ e $V_i$ sono minori dei pesi di tutti gli elementi in $U_{i+1}$ e $V_{i+1}$, per $1 \leq i \leq k-1$.
    \item \textbf{Step 2:} Il processore $P_i$ esegue il merge sequenziale di $U_i$ e $V_i$, per $1 \leq i \leq k$, in parallelo con gli altri processori.
\end{enumerate}

\begin{algorithm}[H]
\caption{Search max-value}\label{alg:parallel_max}
\begin{algorithmic}
\Require Due liste ordinate $A$ e $B$
\Ensure $MaxA$ e $MaxB$
\State Dividi sia $A$ che $B$ uniformemente in $k$ blocchi
\For{\textbf{all} $P_i$, $1 \leq i \leq k$, \textbf{in parallel}}
    \State $MaxA_i \gets a_{i,1}.p$, \quad $MaxB_i \gets b_{i,1}.p$
    \For{$j \gets 2$ \textbf{to} $e$}
        \If{$a_{i,j}.p > MaxA_i$}
            \State $MaxA_i \gets a_{i,j}.p$
        \EndIf
        \If{$b_{i,j}.p > MaxB_i$}
            \State $MaxB_i \gets b_{i,j}.p$
        \EndIf
    \EndFor
\EndFor
\end{algorithmic}
\end{algorithm}

\section{Proposta di progetto}
La proposta concerne l'implementazione di questa metodologia COPA. Prima con OpenMP e poi con MPI, lo scopo è valutare lo speedup ottenibile da quest'implementazione confrontandolo con l'approccio che usa la programmazione dinamica. I punti chiave sono i seguenti:
\begin{enumerate}
    \item Implementare l'algoritmo knapsack con l'approccio DP
    \item Implementare l'algoritmo knapsack sequenziale proposto dall'articolo 
    \item Valutare le performance dei due algoritmi precedenti
    \item Implementare l'algoritmo COPA proposto dall'articolo con OpenMP 
    \item Valutare lo speedup all'aumentare delle unità di esecuzione
    \item Implementare l'algoritmo COPA proposto con MPI 
    \item Valutare lo speedup all'aumentare delle unità di esecuzione
    \item Comparare le performance previste dall'articolo con quelle ottenute sperimentalmente
\end{enumerate}

Il risultato atteso da questo lavoro è di: 1) verificare sperimentalmente che l'algoritmo introduca effettivamente uno speedup. 2) Implementare l'algoritmo con MPI e OpenMP 3) Intodurre se possibile delle modifiche all'algoritmo che permettano di estrarre quali sono gli oggetti scelti. 

Il progetto verrà implementato usando GCC 15 in C++, usando la libreria Boost::MPI, che implementa le funzionalità di MPI sfruttando il sistema di tipizzazione e template di C++, la quale è basata su OpenMPI. 

\bibliographystyle{plain}
\bibliography{biblio}

\end{document}



